\documentclass[10pt,a4paper]{amsart}
\usepackage[margin=19mm]{geometry}
\usepackage{amsmath,amssymb,mathtools}
\usepackage[scr=rsfs,cal=euler]{mathalfa}
\usepackage{xfrac}
\usepackage[unicode,hidelinks,bookmarks=false]{hyperref}
\newcommand{\ie}{i.e.\ }
\newcommand{\cf}{cf.\ }
\newcommand{\resp}{respectively\ }
\newcommand{\Subsection}{\subsection}
\providecommand{\nobreakdash}{\nobreak\hbox{-}}
\setlength{\emergencystretch}{2em}
\multlinegap0pt
\usepackage{bm}
\usepackage{bbm}
\usepackage{dsfont}
\usepackage{xspace}
\usepackage{cancel}
\usepackage{mathtools}
\usepackage{amsthm}
\makeatletter
\@ifundefined{theorem}{\newtheorem{theorem}{Theorem}}{}
\makeatother
\title[The elementary theory of full $n$-branching ordinal trees wi]{The elementary theory of full $n$-branching ordinal trees with successor functions}
\author{Junhong Chen, Yi Zhang}
\date{Source: arXiv:2608.08426v1 (CC BY 4.0)}
\begin{document}
\maketitle

\noindent\emph{Source and licence.} The elementary theory of full $n$-branching ordinal trees with successor functions. Version arXiv:2608.08426v1 (CC BY 4.0), distributed under the Creative Commons Attribution 4.0 International licence, as recorded in the arXiv licence metadata for this exact version and shown on the abstract page https://arxiv.org/abs/2608.08426v1. Full rights notice: licenses/arxiv-2608.08426-CC-BY-4.0.txt.\par\smallskip

\noindent\textit{Editorial scope.} Counted unit: the Theorem whose source label is \texttt{lem:exact-stretching} in the article (source0.tex lines 578--580, source label \texttt{lem:exact-stretching}), delivered together with the article's own complete original proof of it (source0.tex lines 581--623, one \texttt{proof} environment of 43 lines). No mathematics is written, repaired or re-derived: statement and proof are the article's own text line for line. Required context retained uncounted: lem:finite-profile-scheme (lines 208--256); thm:refined-bound (lines 314--316). Every definition, display and label of those lines is retained unchanged. Omitted from this package and not redistributed: 1--207, 257--313, 317--577, 624--1172. Nothing inside a retained excerpt is omitted or altered: every retained body is delivered exactly as the article's lines are written, so no omission receipt of any kind accompanies this package. Citations: the retained excerpts carry no citation command at all, so no bibliography entry of the article is transcribed and no reference is redistributed. Reference adaptation: none was needed and none was made; every reference inside the delivered unit resolves inside it, so no unresolved reference or citation is produced. Printed slips kept literal: no printed word, symbol, hypothesis or number of the counted statement or of its proof was altered. Layout and class: the article's own class is \texttt{amsart} with packages including latexsym, amsfonts, amsmath, amssymb, mathrsfs, tikz-cd, adjustbox, changepage, calc, inputenc, csquotes, braket, bbm, hyperref; the wrapper is the lane's pinned \texttt{amsart} editorial wrapper at 10pt on A4, which restates the theorem-like environments the retained text needs and adds the article's own macro definitions that the retained text needs (none), with extra wrapper packages (bm, bbm, dsfont, xspace, cancel, mathtools). The delivered numbering is the unit's own. Assets: no retained span contains a figure environment, a graphics-inclusion command, a PDF-page inclusion, a table or a TikZ picture; no image file is copied into this package, the package declares no asset dependency and no image file is redistributed with it. Licence: version arXiv:2608.08426v1 of this article is distributed under the Creative Commons Attribution 4.0 International licence, as recorded in the arXiv licence metadata for this exact version and shown on the abstract page https://arxiv.org/abs/2608.08426v1. A full rights notice is delivered at licenses/arxiv-2608.08426-CC-BY-4.0.txt. Characters: every retained excerpt is pure ASCII, so the delivered engine, XeLaTeX with its default Latin Modern text font, reports no missing character.
\begin{theorem}[Finite ordinal-profile scheme]
\label{lem:finite-profile-scheme}
Let
$
  L_m=\{<,P_0,\ldots,P_{m-1}\}$
be a finite relational language with a linear order and unary predicates,
and let $\varphi$ be an $L_m$-sentence.  There are finite data
\[
 \mathcal A_\varphi=(D,Q,I,F_{\mathrm{succ}},F_{\mathrm{lim}},
                      \delta_{\mathrm{next}},\delta_{\mathrm{lim}}),
  Q\subseteq\mathcal P(D),
\]
where $D$ is finite,
\[
 I,F_{\mathrm{succ}}\subseteq Q,\qquad
 F_{\mathrm{lim}}\subseteq\mathcal P(D),\qquad
 \delta_{\mathrm{next}}\subseteq Q^2,\qquad
 \delta_{\mathrm{lim}}\subseteq\mathcal P(D)\times Q,
\]
and a decoding map $\ell:Q\to\{0,1\}^{m}$ with the following property.
For a map $r:\alpha\to Q$ and a nonzero limit ordinal
$\delta\leq\alpha$, define
\[
 \operatorname{LimSupp}(r,\delta)=
 \{e\in D:\exists\beta<\delta\ \forall\xi\,
       (\beta<\xi<\delta\longrightarrow e\in r(\xi))\}.
\]
For every nonzero ordinal $\alpha$ and every $L_m$-structure
\[
 \mathfrak M=(\alpha,<,P_0^{\mathfrak M},\ldots,
                         P_{m-1}^{\mathfrak M}),
\]
we have $\mathfrak M\models\varphi$ if and only if there is a map
$r:\alpha\to Q$ such that
\begin{enumerate}
  \item $r(0)\in I$, and
  $\ell(r(\xi))_i=1$ if and only if
  $\xi\in P_i^{\mathfrak M}$, for every $\xi<\alpha$ and $i<m$;
  \item $(r(\xi),r(\xi+1))\in\delta_{\mathrm{next}}$ whenever
  $\xi+1<\alpha$;
  \item $(\operatorname{LimSupp}(r,\delta),r(\delta))\in
  \delta_{\mathrm{lim}}$ for every nonzero limit
  $\delta<\alpha$;
  \item if $\alpha=\beta+1$, then $r(\beta)\in F_{\mathrm{succ}}$,
  and if $\alpha$ is a limit ordinal, then
  $\operatorname{LimSupp}(r,\alpha)\in F_{\mathrm{lim}}$.
\end{enumerate}
The data are finite and effectively obtainable from $\varphi$.
\end{theorem}

\begin{theorem}[Refined bound]\label{thm:refined-bound}
    Let $q\geq2$, $k=|\rho_{q-1}(\mathfrak A)|$. For every $\mathfrak A$, there is a $\mathfrak B$ that is $q$-equivalent to $\mathfrak A$ and $\operatorname{otp}(B)\leq\omega^k$.
\end{theorem}

\begin{theorem}[Exact stretching]\label{lem:exact-stretching}
    Let $\kappa$ and $\lambda$ be regular uncountable cardinals. For every fixed $q<\omega$ and every FCO structure $\mathfrak{A}=(\lambda,<,C_0,\dots,C_{d-1})$, there is an FCO structure $\mathfrak{B}=(\kappa,<,D_0,\dots,D_{d-1})$ such that $\mathfrak{A}\equiv_q\mathfrak{B}$.
\end{theorem}

\begin{proof}

Apply Theorem~\ref{lem:finite-profile-scheme} to $[[\mathfrak A]]^q$, and use its notation
$D,Q,I,\delta_{\mathrm{next}},\delta_{\mathrm{lim}},F_{\mathrm{lim}}$, and
$\ell$.  Since $\mathfrak A\models\varphi$, where $\varphi=[[(\mathfrak A)]]^q$, there is an accepting profile assignment $r\colon\lambda\to Q$.
Let
$
Y=\operatorname{LimSupp}(r,\lambda).
$
Then $Y\in F_{\mathrm{lim}}$.

 For $e\notin Y$, put $S_e=\{\xi<\lambda:e\notin r(\xi)\}.$
Each $S_e$ is unbounded in $\lambda$.  Since $\lambda$ is regular
uncountable, the set
$
  \operatorname{Acc}(S_e)=
  \{\delta<\lambda:S_e\cap\delta\text{ is unbounded in }\delta\}$
  is club in $\lambda$.  Hence the finite intersection
   $\cap_{e\notin Y}\operatorname{Acc}(S_e)$
Choose $\delta_0<\delta_1$ in this club above a stage after which every $e\in Y$ belongs to every $r(\xi)$. Then
$
  \operatorname{LimSupp}(r,\delta_0)=
  \operatorname{LimSupp}(r,\delta_1)=Y.
 $

The profile
condition at $\delta_0$ gives$  (Y,r(\delta_0))\in\delta_{\mathrm{lim}}.$

Let $U$ be the colored ordinal
$r\mathbin{\upharpoonright}[0,\delta_0)$, and let $V$ be the reindexed
colored interval
$r\mathbin{\upharpoonright}[\delta_0,\delta_1)$.  Both
have no last point.

We record the finite profiles of $U$ and $V$ in the finite profile language. By the bound Theorem ~\ref{thm:refined-bound}, we obtain countable copies $U^*$ and $V^*$ satisfying the corresponding profile sentences.

Now, we can consider
\[
  \mathfrak B= \mathfrak U^*+\sum_{\xi<\kappa}\mathfrak V^*_{\xi}.
\]
The profile obtained by concatenating $U^*$ with $\kappa$ copies of $V^*$ is accepted by the same automaton and decodes to $\mathfrak B$. Hence $\mathfrak B\models[[(\mathfrak A)]]^q$, and therefore $\mathfrak A\equiv_q\mathfrak B$.

\end{proof}

\end{document}
